% =============================================================================
% Materiał i metody. Plik edytowany ręcznie -- skrypt analiza_odpowiedzi_naukowa.py go NIE nadpisuje.
% Liczby: definicje.tex; tabele: \TabDir (CSV); rysunki: \FigDir; literatura: references.bib.
% Analizy uzupełniające (zmiana na poziomie pracowników, analiza pytań) opisano w suplement.tex.
% =============================================================================
\section{Materiał i metody}

\subsection{Środowisko i program szkoleniowy}
Program prowadzono w banku. Jego treść stanowiły pytania ogólne pogrupowane w cztery kategorie:
wymogi regulacyjne, produkty bankowe, umiejętności miękkie oraz moduł dla kadry kierowniczej.
Każde pytanie było pytaniem wielokrotnego wyboru z trzema do pięciu opcjami odpowiedzi, z których
dokładnie jedna była poprawna, i zawierało wyjaśnienie poprawnej odpowiedzi. Aplikacja szkoleniowa
uruchamiała się automatycznie na ekranach komputerów pracowników i wyświetlała od jednego do trzech
pytań dziennie w godzinach 10:00--14:00; po każdej odpowiedzi pokazywano poprawną odpowiedź wraz
z wyjaśnieniem. To samo pytanie mogło zostać ponownie zadane temu samemu pracownikowi nie wcześniej
niż siedem dni po poprzedniej odpowiedzi, a harmonogram powtórek nie zależał od tego, czy odpowiedź
była poprawna. Na czas przeglądu i aktualizacji pytań program był zawieszany.
O udziale w programie decydował Departament Digitalizacji i Procesów banku, który był również autorem
pytań. Pytania prezentowano w modułach tematycznych (mediana \OptionsMedian{} zaobserwowanych opcji
odpowiedzi na pytanie). W analizowanych danych pracownicy odpowiadali na medianę \DailyMedian{}
pytań w dniu aktywności, a to samo pytanie było ponownie zadawane temu samemu pracownikowi po
medianie \GapMedian{} dni (IQR \GapQOne--\GapQThree).

\subsection{Dane}
Dane przekazano w formie listy odpowiedzi zawierającej moduł, pytanie, wybraną opcję, kod
odpowiadającego, znacznik czasu oraz poprawność odpowiedzi (wartość logiczna). Wykluczono
\NExcludedItems{} pytania, w których zarejestrowano tylko jedną opcję odpowiedzi
(\NExcludedItemResp{} odpowiedzi), ponieważ nie niosą one informacji o błędach. Zbiór analityczny
obejmował \NResponses{} odpowiedzi (tab.~\ref{tab:zbior_danych}). Dane osób odpowiadających
przekazano w formie metadanych zawierających płeć pracownika oraz przedział wiekowy (co 10 lat) oraz medianę wieku.
% SPRAWDŹ: people.csv - nie bedzie dostarczany i brany pod uwage.

\fillin{Podstawa prawna przetwarzania danych spseudonimizowanych (wiek, płeć)}

\subsection{Definicje}
\emph{Sekwencja} to uporządkowana chronologicznie seria prób jednego pracownika dotyczących jednego
pytania; \emph{numer próby} to pozycja odpowiedzi w jej sekwencji (przy jednakowym czasie o kolejności
decydował identyfikator rekordu). \emph{Poprawność} to odsetek poprawnych odpowiedzi. \emph{Odstęp}
to czas w dniach od poprzedniej próby w tej samej sekwencji. \emph{Staż} to czas od pierwszej
odpowiedzi pracownika w programie.

\subsection{Analiza opisowa}
Odsetki poprawnych odpowiedzi podano z 95-procentowymi przedziałami ufności (PU) z bootstrapu
klastrowego, w którym losowano pracowników ze zwracaniem (\BootReps{} replikacji); procedura ta
uwzględnia zależność odpowiedzi udzielanych przez tę samą osobę. Ponieważ osoby kontynuujące do
kolejnych prób mogą różnić się od tych, które przestają, zagregowane krzywe uczenia mogą być
zniekształcone przez ubywanie obserwacji \citep{murray2013,goutte2018}; krzywą obliczono więc
dodatkowo dla zbalansowanej kohorty sekwencji liczących co najmniej \CohortK{} prób. Zależność odsetka
błędów od numeru próby opisano funkcją potęgową \citep{newell1981}. Ponieważ harmonogram powtórek nie
był adaptacyjny, sprawdzono, czy ponowne zadanie pytania zależało od wyniku poprzedniej próby.

\subsection{Modele statystyczne}
Niech $Y_{ijt}$ oznacza poprawność ($1$ albo $0$) odpowiedzi pracownika $i$ na pytanie $j$ w próbie
$t$. Zastosowano trzy modele regresji logistycznej o wspólnej postaci
\begin{equation*}
\operatorname{logit} P(Y_{ijt}=1) = \alpha_i + \gamma_j + \mathbf{x}_{ijt}^{\top}\boldsymbol{\beta},
\end{equation*}
gdzie $\alpha_i$ i $\gamma_j$ są efektami stałymi pracownika i pytania. Efekty stałe pochłaniają
stałe różnice poziomu wiedzy między osobami i trudności między pytaniami, dlatego współczynniki
$\boldsymbol{\beta}$ opisują zmiany w obrębie tej samej osoby i tego samego pytania. Modele różniły
się zbiorem odpowiedzi i wektorem predyktorów $\mathbf{x}_{ijt}$:
\begin{description}
\item[M1, model uczenia (PB1).] Wszystkie odpowiedzi; jedynym predyktorem jest numer próby,
  $\beta_1 \log_2 t$. Wartość $e^{\beta_1}$ jest ilorazem szans (OR) poprawnej odpowiedzi przy
  podwojeniu numeru próby (np.\ z 1 do 2 albo z 4 do 8).
\item[M2, model retencji (PB2).] Wyłącznie próby powtórne ($t \geq 2$);
  $\beta_1 \log_2 t + \beta_2 Y_{ij,t-1} + \beta_3 \log_2 \Delta_{ijt} + \beta_4 S_{it}$, gdzie
  $Y_{ij,t-1}$ to wynik poprzedniej próby, $\Delta_{ijt}$ odstęp od niej w dniach (ograniczony od dołu
  do 1 dnia), a $S_{it}$ staż w latach. Konstrukcją model odpowiada modelom \emph{performance factors
  analysis} i \emph{half-life regression} stosowanym do logów uczących się \citep{pavlik2009,settles2016}.
\item[M3, model nowych pytań.] Wyłącznie pierwsze próby ($t = 1$); jedynym predyktorem jest staż,
  $\beta_4 S_{it}$. Przy efektach stałych pracownika współczynnik porównuje pierwsze kontakty tej samej
  osoby z kolejnymi pytaniami, napotkanymi przy różnym stażu.
\end{description}
Pracowników i pytania, w których wszystkie odpowiedzi były poprawne albo wszystkie błędne, usuwano
iteracyjnie przed estymacją, ponieważ ich efekty stałe nie są identyfikowalne. Modele szacowano metodą
największej wiarygodności, a błędy standardowe klastrowano względem pracowników.

Odporność oszacowania $e^{\beta_1}$ z modelu M1 sprawdzono w analizach wrażliwości: w zbalansowanej
kohorcie, po obcięciu sekwencji do 10 prób, w warstwach długości sekwencji (kontrola obciążenia typu
\emph{pattern-mixture}), po wykluczeniu pracowników z mniej niż 50 odpowiedziami oraz po dodaniu stażu
jako predyktora. Efekt uśredniony populacyjnie oszacowano ponadto metodą uogólnionych równań
estymujących (GEE) z wymienną macierzą korelacji roboczej \citep{liang1986}, a w modelu mieszanym
efekty stałe zastąpiono skrzyżowanymi losowymi wyrazami wolnymi pracowników i pytań \citep{baayen2008}.
Z modelu mieszanego wyznaczono współczynniki korelacji wewnątrzklasowej (ICC) w skali latentnej oraz
brzegowy i warunkowy współczynnik $R^2$ \citep{nakagawa2013}.

\subsection{Przejścia i powtarzane błędy}
Przejścia między kolejnymi próbami tego samego pytania zestawiono w tabeli kontyngencji. Dla dwóch
następujących po sobie błędów porównano odsetek ponownych wyborów tego samego dystraktora z odsetkiem
oczekiwanym przy losowym wyborze spośród zaobserwowanych dystraktorów oraz przy wyborze proporcjonalnym
do popularności każdego dystraktora w danym pytaniu.

Analizy uzupełniające, czyli zmianę poprawności na poziomie pracowników oraz analizę pytań metodami
klasycznej teorii testu, opisano w materiale uzupełniającym. Obliczenia wykonano w języku Python
\VerPython{} z użyciem bibliotek pandas \VerPandas, NumPy \VerNumpy, SciPy \VerScipy, statsmodels
\VerStatsmodels{} i Matplotlib \VerMatplotlib{}.
%  (ziarno generatora liczb losowych \Seed).
